% ID: SEL-D0058-OB-E02
% ESTADO: validada. S058.
\textbf{a)} Las ramas son continuas en sus tramos y el dominio es todo $\mathbb R$, pues en la rama racional $x>2$. En $x=2$, igualamos $f(2)=2$ y el límite derecho $4-a/2$, de donde $\boxed{a=4}$. Las derivadas de los tramos son $2x-3$ y $4/x^2$. Las derivadas laterales en $2$ coinciden y valen $1$, de modo que con $a=4$ la función es derivable en todo $\mathbb R$.

\textbf{b)} Con $a=1$ no hay asíntotas verticales. La rama racional no se usa cerca de $0$ y en $2$ los límites laterales son finitos, $2$ y $7/2$. Al tender $x$ a $+\infty$, $f(x)\to4$, por lo que $\boxed{y=4}$ es asíntota horizontal por la derecha. A la izquierda la parábola tiende a $+\infty$ y no hay asíntota horizontal.
